\newcommand{\pendcorto}{\textcolor{red!70!black}{\textbf{[PEND.]}}}
% =====================================================================
% MEJORADO (27 oct 2026): pruebas de seguridad (reglas y funciones), umbral de reconocimiento,
% Android 7.0 (API 24) y tabla de desviaciones.
\section{Prueba final}
\label{sec:des-pruebas}

La última fase de Mobile-D verifica que el producto cumpla sus
requerimientos. Las pruebas se organizaron en seis niveles, de los más
automáticos a los que requieren al usuario (Tabla~\ref{tab:estrategia}).

\begin{table}[H]
\centering
\caption{Estrategia de pruebas}
\label{tab:estrategia}
\begin{tabularx}{\textwidth}{|L{3.5cm}|Y|L{3.4cm}|}
\hline
\rowcolor{encabezado} \textbf{Nivel} & \textbf{Objeto de la prueba} & \textbf{Técnica} \\
\hline
Verificación estática & Tipos y calidad de todo el código. & Compilador y
analizadores (sección~\ref{sec:des-estabilizacion}) \\ \hline
Lógica de negocio & Funciones puras: generador de rutinas, IMC, estadísticas
de entrenamiento, porciones de alimentos y traducción de consultas. &
Pruebas unitarias con valores límite \\ \hline
Reglas de seguridad y funciones & Aislamiento entre usuarios, validación de
datos y control de acceso en el servidor. & Escenarios permitidos y denegados
(emulador de Firebase) \\ \hline
Funcionales & Cada requerimiento funcional en la aplicación instalada. &
Caja negra, casos de prueba \\ \hline
Usabilidad & Facilidad de uso de las tareas principales. & Tareas con
usuarios \\ \hline
Aceptación & Cumplimiento de cada objetivo específico. & Trazabilidad
objetivo--requerimiento--prueba \\
\hline
\end{tabularx}
\end{table}

\subsection{Verificación técnica}

\subsubsection{Pruebas de la lógica de negocio}
\label{subsec:pruebas-logica}

Las funciones de las reglas de negocio no dependen de la red ni de la base de
datos, por lo que se ejercitaron directamente con Node.js~22 sobre el código
fuente del proyecto, sin modificar su lógica. Las funciones del controlador de
nutrición se copiaron textualmente a un archivo auxiliar para aislarlas del SDK
de Firebase. La ejecución se realizó el 14 de octubre de 2026.

\paragraph{Generador de rutinas.} Se evaluaron las doce combinaciones posibles
de objetivo y nivel. La Tabla~\ref{tab:prueba-generador} muestra los
resultados obtenidos, que coinciden en todos los casos con las reglas de las
Tablas~\ref{tab:reglas-frecuencia} a~\ref{tab:reglas-division}.

\begin{table}[H]
\centering
\caption{Resultados del generador para las doce combinaciones de entrada}
\label{tab:prueba-generador}
\small
\begin{tabular}{|L{2.5cm}|L{2.2cm}|C{0.9cm}|C{1.9cm}|L{4.4cm}|C{1.1cm}|}
\hline
\rowcolor{encabezado} \textbf{Objetivo} & \textbf{Nivel} & \textbf{Días} & \textbf{Series × Rep.
(desc.)} & \textbf{División obtenida} & \textbf{Ejerc.} \\
\hline
Ganar masa & Principiante & 3 & 4×9 (90 s) & Emp. -- Tir. -- Pie. & 15 \\ \hline
Ganar masa & Intermedio & 4 & 4×9 (90 s) & Emp. -- Tir. -- Pie. -- C.C. & 21 \\ \hline
Ganar masa & Avanzado & 5 & 4×9 (90 s) & Emp. -- Tir. -- Pie. -- Emp. -- Tir. & 25 \\ \hline
Mantenimiento & Principiante & 3 & 3×12 (60 s) & Emp. -- Tir. -- Pie. & 15 \\ \hline
Mantenimiento & Intermedio & 4 & 3×12 (60 s) & Emp. -- Tir. -- Pie. -- C.C. & 21 \\ \hline
Mantenimiento & Avanzado & 5 & 3×12 (60 s) & Emp. -- Tir. -- Pie. -- Emp. -- Tir. & 25 \\ \hline
Perder peso & Principiante & 3 & 3×15 (45 s) & Emp. -- Car. -- Pie. & 14 \\ \hline
Perder peso & Intermedio & 4 & 3×15 (45 s) & Emp. -- Tir. -- Car. -- Pie. & 19 \\ \hline
Perder peso & Avanzado & 5 & 3×15 (45 s) & Emp. -- Tir. -- Pie. -- Car. -- C.C. & 25 \\ \hline
Resistencia & Principiante & 3 & 3×18 (30 s) & Emp. -- Car. -- Pie. & 14 \\ \hline
Resistencia & Intermedio & 4 & 3×18 (30 s) & Emp. -- Tir. -- Car. -- Pie. & 19 \\ \hline
Resistencia & Avanzado & 5 & 3×18 (30 s) & Emp. -- Tir. -- Pie. -- Car. -- C.C. & 25 \\
\hline
\end{tabular}
\fuente{Ejecución de \archivo{generarPlanHeuristico} (comportamiento base, sin
edad, IMC ni días disponibles). Emp.: empuje; Tir.:
tirón; Pie.: pierna; Car.: cardio; C.C.: cuerpo completo; Ejerc.: total de
ejercicios por semana.}
\end{table}

\paragraph{Cálculo del IMC.} Se probaron valores en los límites de cada
categoría de la OMS (Tabla~\ref{tab:prueba-imc}). Todos los casos produjeron
el resultado esperado.

\begin{table}[H]
\centering
\caption{Pruebas de valores límite del cálculo y la clasificación del IMC}
\label{tab:prueba-imc}
\small
\begin{tabular}{|C{2.2cm}|C{2.4cm}|C{2.2cm}|L{3.3cm}|C{2cm}|}
\hline
\rowcolor{encabezado} \textbf{Peso (kg)} & \textbf{Estatura (cm)} & \textbf{IMC} &
\textbf{Categoría obtenida} & \textbf{Resultado} \\
\hline
50.0 & 170 & 17.3 & Bajo peso & Correcto \\ \hline
53.4 & 170 & 18.5 & Normal & Correcto \\ \hline
70.0 & 175 & 22.9 & Normal & Correcto \\ \hline
72.2 & 170 & 25.0 & Sobrepeso & Correcto \\ \hline
86.7 & 170 & 30.0 & Obesidad & Correcto \\ \hline
100.0 & 170 & 34.6 & Obesidad & Correcto \\ \hline
\multicolumn{3}{|c|}{Clasificación de 18.4 / 18.5} & Bajo peso / Normal & Correcto \\ \hline
\multicolumn{3}{|c|}{Clasificación de 24.9 / 25.0} & Normal / Sobrepeso & Correcto \\ \hline
\multicolumn{3}{|c|}{Clasificación de 29.9 / 30.0} & Sobrepeso / Obesidad & Correcto \\
\hline
\end{tabular}
\fuente{Ejecución de \archivo{calcularIMC} y \archivo{clasificarIMC}.}
\end{table}

\paragraph{Reparto, entrenamientos, porciones y traducción.} La
Tabla~\ref{tab:prueba-otras} reúne las demás pruebas de funciones puras. Para
las estadísticas de entrenamiento se fijó la fecha actual en el miércoles 14 de
octubre de 2026.

\begin{table}[H]
\centering
\caption{Otras pruebas de la lógica de negocio}
\label{tab:prueba-otras}
\footnotesize
\begin{tabularx}{\textwidth}{|L{1.1cm}|L{4.1cm}|Y|Y|}
\hline
\rowcolor{encabezado} \textbf{ID} & \textbf{Función} & \textbf{Entrada} & \textbf{Resultado obtenido} \\
\hline
PL-01 & \archivo{repartirCantidad} & 5 ejercicios en 3 categorías & 2 -- 2 -- 1 \\ \hline
PL-02 & \archivo{repartirCantidad} & 6 ejercicios en 3 categorías & 2 -- 2 -- 2 \\ \hline
PL-03 & \archivo{repartirCantidad} & 4 ejercicios en 2 categorías & 2 -- 2 \\ \hline
PL-04 & Estadísticas de entrenamiento & Días 12, 13 y 14 de octubre (hoy
incluido) & Racha 3; semana 3; total 3 \\ \hline
PL-05 & Estadísticas de entrenamiento & Días 12 y 13 (hoy sin completar) &
Racha 2; semana 2 \\ \hline
PL-06 & Estadísticas de entrenamiento & Días 10 y 12 (falta el día 13) &
Racha 0; semana 1 \\ \hline
PL-07 & Estadísticas de entrenamiento & Diez días consecutivos hasta hoy &
Racha 10; semana 3; total 10 \\ \hline
PL-08 & Estadísticas de entrenamiento & Sin registros & Racha 0; semana 0;
total 0 \\ \hline
PL-09 & \archivo{calcularPorcion} & Arroz (130~kcal, 2.7~g P, 28~g C, 0.3~g
G por 100~g), 200~g & 260~kcal; 5.4~g P; 56~g C; 0.6~g G \\ \hline
PL-10 & \archivo{calcularPorcion} & Pollo (165~kcal, 31~g P por 100~g), 150~g
& 248~kcal; 46.5~g P; 0~g C; 5.4~g G \\ \hline
PL-11 & \archivo{calcularPorcion} & Cualquier alimento, 0~g & Todos los
valores en 0 \\ \hline
PL-12 & \archivo{calcularResumenDiario} & Dos alimentos de hoy (arroz y pollo),
uno de otra fecha y un cálculo de IMC & Solo suma los dos de hoy: 508~kcal,
quedan 1\,492; 25.4\,\% de la meta de referencia (2\,000~kcal) \\ \hline
PL-13 & \archivo{calcularResumenDiario} & Un alimento de 2\,600~kcal & Porcentaje
limitado a 100\,\%; restantes 0 \\ \hline
PL-14 & \archivo{traducirConsulta} & ``pechuga de pollo'', ``Plátano'',
``leche descremada'' & ``chicken breast'', ``banana'', ``milk nonfat'' \\ \hline
PL-15 & \archivo{traducirConsulta} & ``chicken'' y ``a'' & ``chicken'' sin
cambio; ``a'' no genera consulta \\
\hline
\end{tabularx}
\end{table}

\subsubsection{Pruebas unitarias automáticas}
\label{sec:pruebas-unitarias}

Las pruebas anteriores se ejecutaron una vez, con \emph{scripts} auxiliares.
Para que las reglas de negocio puedan verificarse de forma repetible después de
cada cambio, se escribieron 20 archivos de prueba de Jest
(Tabla~\ref{tab:suites-jest}), que se ejecutan con \archivo{npm test}. Las
funciones probadas son puras, por lo que las pruebas no requieren Firebase ni
un dispositivo. Entre los casos cubiertos están:

\begin{itemize}
    \item \textbf{IMC:} los límites exactos de cada categoría de la OMS
    (18,4 es bajo peso; 18,5 es normal; 24,9 es normal; 25 es sobrepeso; 30 es
    obesidad).
    \item \textbf{Metas nutricionales:} que un hombre de 25 años, 75~kg y
    178~cm, que entrena cuatro días para ganar masa, obtenga una tasa basal de
    1\,743~kcal y una meta de 3\,001~kcal; que perder peso aplique un déficit de
    500~kcal; que nunca se baje del piso calórico (1\,200~kcal para mujeres y
    1\,500~kcal para hombres); y que los macronutrientes no sean negativos y
    reconstruyan las calorías con una tolerancia de 10~kcal.
    \item \textbf{Generador de rutinas:} que el orden de las seis plantillas y
    de los cuatro esquemas coincida con el del \emph{script} de entrenamiento
    del modelo (si cambiara, el modelo dejaría de ser válido); que IMC $\geq$ 30
    active el énfasis en cardio; que los mayores de 60 años no reciban
    hipertrofia ni esquema metabólico y los menores de 16 no reciban
    hipertrofia; y que cualquier combinación de perfil produzca clases
    válidas.
    \item \textbf{Vector de entrada del modelo:} el orden de las diez
    características, la codificación \emph{one-hot} del objetivo y los valores
    por defecto cuando faltan datos del perfil.
    \item \textbf{Estadísticas y sesión de entrenamiento:} racha, semana y
    total; creación de la sesión, series completadas y descansos.
    \item \textbf{Reconocimiento y realidad aumentada:} que las cajas de las
    candidatas alternativas no se mezclen con las de la máquina principal, que
    el objetivo que recibe la vista de ARCore tome la caja de la candidata
    confirmada, y que el esquema de la respuesta rechace cajas inválidas.
    \item \textbf{Otros componentes:} carga diferida de TensorFlow Lite,
    esquemas de los datos persistidos, actualización parcial del progreso,
    persistencia de la sesión de entrenamiento, demostraciones locales por
    máquina y verificación de las variables de entorno.
\end{itemize}

\subsubsection{Pruebas de las reglas de seguridad y de las funciones del servidor}
\label{sec:pruebas-seguridad}

Las reglas de seguridad y las funciones del servidor constituyen la barrera
frente a un acceso malintencionado, porque las credenciales de Firebase están
contenidas en la aplicación y no son secretas; App Check añade una verificación
complementaria de la procedencia de las solicitudes. La Tabla~\ref{tab:prueba-reglas} define los escenarios evaluados
y el resultado que exigen las reglas del Anexo~\ref{anx:reglas}. Los
escenarios se automatizaron con el emulador de Firebase (comando
\archivo{npm run test:security}), autenticando como un usuario $A$ y operando
sobre datos propios o de otro usuario $B$: el archivo
\archivo{tests/security/rules.test.mjs} declara 11 pruebas sobre Firestore y
Storage, y \archivo{functions/test/security.integration.mjs} declara 4 pruebas
de integración de las funciones. Las pruebas de los esquemas del servidor
(\archivo{npm run test:backend}), que incluyen la interpretación de la
respuesta del reconocimiento (clases, equivalencias, cajas inválidas y confianza
fuera de rango), se ejecutaron el 3 de noviembre de 2026 y aprobaron 22 de 22
casos. Los 19 escenarios de seguridad de la
Tabla~\ref{tab:prueba-reglas} también se completaron satisfactoriamente: los
resultados obtenidos coincidieron con los esperados.

\begin{table}[H]
\centering
\caption{Escenarios de prueba de las reglas de seguridad}
\label{tab:prueba-reglas}
\footnotesize
\begin{tabularx}{\textwidth}{|L{1.3cm}|Y|C{1.8cm}|C{1.9cm}|}
\hline
\rowcolor{encabezado} \textbf{ID} & \textbf{Escenario} & \textbf{Esperado} & \textbf{Obtenido} \\
\hline
SEG-01 & Usuario no autenticado lee \archivo{maquinas/\{id\}} & Denegado & Denegado \\ \hline
SEG-02 & $A$ lee \archivo{maquinas/\{id\}} & Permitido & Permitido \\ \hline
SEG-03 & $A$ escribe en \archivo{maquinas/\{id\}} & Denegado & Denegado \\ \hline
SEG-04 & $A$ lee \archivo{usuarios/A/rutinas} & Permitido & Permitido \\ \hline
SEG-05 & $A$ lee \archivo{usuarios/B/rutinas} & Denegado & Denegado \\ \hline
SEG-06 & $A$ envía a \archivo{guardarRutina} una rutina con
\archivo{diasPorSemana = 8} & Rechazado & Rechazado \\ \hline
SEG-07 & $A$ envía a \archivo{guardarRutina} una rutina con
\archivo{usuarioId = B} (el propietario se toma de la identidad autenticada) & Rechazado & Rechazado \\ \hline
SEG-08 & $A$ modifica el campo \archivo{email} de su perfil & Denegado & Denegado \\ \hline
SEG-09 & $A$ crea su perfil con \archivo{edad = 10} & Denegado & Denegado \\ \hline
SEG-10 & $A$ registra progreso con \archivo{peso = 0} & Denegado & Denegado \\ \hline
SEG-11 & $A$ registra ingesta con \archivo{calorias = -100} & Denegado & Denegado \\ \hline
SEG-12 & $A$ crea un entrenamiento cuyo \archivo{fecha} difiere del
identificador del documento & Denegado & Denegado \\ \hline
SEG-13 & $A$ crea un entrenamiento con \archivo{dia = 8} & Denegado & Denegado \\ \hline
SEG-14 & $A$ crea su perfil con \archivo{creadoEn} distinto de la hora del
servidor & Denegado & Denegado \\ \hline
SEG-15 & $A$ sube un archivo a Storage & Denegado & Denegado \\ \hline
SEG-16 & $A$ crea directamente un documento en \archivo{usuarios/A/rutinas},
sin usar la función & Denegado & Denegado \\ \hline
SEG-17 & Llamada a \archivo{guardarRutina}, \archivo{buscarAlimentos} o
\archivo{reconocerMaquina} sin autenticación & Rechazado & Rechazado \\ \hline
SEG-18 & Llamada a las funciones con un usuario anónimo & Rechazado & Rechazado \\ \hline
SEG-19 & Solicitudes concurrentes que superan la cuota por usuario & Limitadas a la cuota & Limitadas a la cuota \\
\hline
\end{tabularx}
\end{table}

\subsubsection{Evaluación de los modelos de inteligencia artificial}
\label{sec:pruebas-ia}

La evaluación distingue el clasificador local, el detector en la nube y el
recomendador de rutinas. Se presentan sus resultados, las condiciones de
evaluación y las limitaciones. Los resultados proceden de los registros
citados en cada tabla; no equivalen a una validación con usuarios.

\paragraph{Evaluación del clasificador local de máquinas}

\paragraph{Evaluación.} Se evaluó con validación cruzada de cinco pliegues y
cuidando que las versiones aumentadas de una misma foto nunca estuvieran en
entrenamiento y validación a la vez. La Tabla~\ref{tab:metricas-reconocimiento}
resume los resultados con imágenes limpias y con el escenario de campo
simulado.

\begin{table}[H]
\centering
\caption{Exactitud del reconocimiento de máquinas (validación cruzada, 35 clases)}
\label{tab:metricas-reconocimiento}
\small
\begin{tabular}{|l|C{2.6cm}|C{2.6cm}|}
\hline
\rowcolor{encabezado} \textbf{Escenario} & \textbf{Acierto top-1} & \textbf{Acierto top-3} \\
\hline
Imágenes limpias del catálogo & 43,4\,\% & 65,9\,\% \\ \hline
Campo simulado (fondos y cámara) & 38,5\,\% & 61,4\,\% \\
\hline
\end{tabular}
\fuente{Archivo \archivo{ml/results/recognition/recognition\_metrics.json}. El
azar acertaría el 2,9\,\% (1 de 35).}
\end{table}

El umbral de confianza se calibró sobre esas mismas predicciones
(Tabla~\ref{tab:umbral-reconocimiento}): a mayor umbral, más seguro es el
reconocimiento automático, pero menos veces se usa. Se eligió 0,8 porque,
cuando la aplicación acepta la máquina sin preguntar, acierta
aproximadamente 85 de cada 100 veces; en las demás ocasiones el usuario
confirma una sugerencia ya preseleccionada.

\begin{table}[H]
\centering
\caption{Calibración del umbral de confianza}
\label{tab:umbral-reconocimiento}
\small
\begin{tabular}{|C{2.2cm}|C{3.2cm}|C{3.2cm}|}
\hline
\rowcolor{encabezado} \textbf{Umbral} & \textbf{Cobertura} & \textbf{Precisión} \\
\hline
0,5 & 44,4\,\% & 60,5\,\% \\ \hline
0,7 & 23,6\,\% & 75,1\,\% \\ \hline
\textbf{0,8} & \textbf{15,8\,\%} & \textbf{85,4\,\%} \\ \hline
0,9 & 9,2\,\% & 92,2\,\% \\
\hline
\end{tabular}
\fuente{Archivo \archivo{ml/results/recognition/recognition\_metrics.json}.
Cobertura: proporción de fotos que superan el umbral; precisión: aciertos entre
ellas.}
\end{table}

\paragraph{Limitaciones.} El resultado es modesto y debe leerse como tal. Con
solo 399 fotografías de catálogo (de una a 25 por máquina) el modelo distingue
bien las máquinas visualmente más características, como la caminadora (93\,\%
de acierto en el campo simulado) o la bicicleta estática (67\,\%), pero casi no
reconoce otras: curl de bíceps, pantorrillas, torso rotatorio, elevaciones
laterales y remo en T se acercan al 0\,\%, y dos máquinas con una sola foto no
pudieron evaluarse. Además, las fotografías de catálogo no son imágenes tomadas
en un gimnasio real. Por ello la máquina se acepta directamente solo cuando la confianza alcanza el
umbral de 0,8, con una cobertura del 15,8\,\% de las fotografías y una
precisión del 85,4\,\% en la validación cruzada; en los demás casos el
reconocimiento se presenta como una \emph{sugerencia} que el usuario confirma o
corrige. La mejora esperada proviene de
más fotografías propias o autorizadas, tomadas en gimnasios reales
(sección~\ref{sec:recomendaciones}).



\paragraph{Evolución y evaluación del detector en la nube}

\paragraph{Primera versión del detector (YOLO11x).} Esta versión parte de
YOLO11x, el tamaño extragrande de la generación YOLO11, con los pesos
preentrenados que distribuye Ultralytics (Khanam \& Hussain, 2024; Jocher
et~al., 2023). Se ajustó en Google Colab con un conjunto anotado en Roboflow
(Roboflow, Inc., s.~f.) y exportado en formato YOLO. El conjunto tiene 458
imágenes y 625 cajas de 42 clases, repartidas en 327 imágenes de entrenamiento,
88 de validación y 43 de prueba. Se entrenó durante 15 épocas con imágenes de
640$\times$640 píxeles y lotes de 16, con los aumentos de datos por defecto de
Ultralytics (mosaico, cambios de color y volteo horizontal). La evaluación se
hizo sobre la partición de prueba, en CPU, tras comprobar con \emph{hashes} que
ninguna imagen de prueba fuera idéntica a una de entrenamiento o validación
(Tabla~\ref{tab:metricas-detector}).

\begin{table}[H]
\centering
\caption{Métricas del detector de máquinas en la partición de prueba}
\label{tab:metricas-detector}
\small
\begin{tabular}{|l|C{2.6cm}|}
\hline
\rowcolor{encabezado} \textbf{Métrica} & \textbf{Valor} \\
\hline
Precisión & 79,8\,\% \\ \hline
Exhaustividad (\emph{recall}) & 29,9\,\% \\ \hline
mAP50 & 39,8\,\% \\ \hline
mAP50--95 & 37,6\,\% \\
\hline
\end{tabular}
\fuente{Archivo \archivo{ml/detection/runs/best42-evaluation/results.json},
generado con \archivo{evaluate\_best42.py}. Partición de prueba: 43 imágenes,
65 cajas, 24 clases con ejemplos; entrada de 640 píxeles; evaluación en CPU de
un computador, no en un teléfono. Definición de las métricas en Padilla
et~al. (2021).}
\end{table}

Estos resultados deben leerse con cautela por cuatro razones. La primera es que
la partición de prueba es mínima: 43 imágenes con 65 cajas, de las cuales 24
pertenecen a una sola clase (\texttt{hip\_adduction\_machine}) y 18 de las 42
clases no tienen ningún ejemplo. La segunda es el desbalance del entrenamiento:
esa misma clase aporta 91 cajas, la mediana es de 8 cajas por clase, y cuatro
clases (extensión de cuádriceps, poleas de doble cable, rack de sentadilla y
remo alto) no tienen ninguna, por lo que el detector no puede reconocerlas. La
tercera es que varias imágenes parecen series de una misma máquina vista desde
ángulos cercanos; aunque no haya duplicados exactos, la prueba puede resultar
más fácil que una situación real en un gimnasio. La cuarta es que una
exhaustividad del 29,9\,\% significa que, en el punto de operación que reporta
Ultralytics, el detector encuentra menos de un tercio de las cajas reales. Por
todo ello las cifras muestran que el procedimiento funciona de extremo a
extremo, no que el detector sea fiable en un gimnasio.

El proyecto incluye además un \emph{script} para exportar la primera versión del detector al formato
LiteRT, el sucesor de TensorFlow Lite, y ejecutarlo en el propio teléfono
(\archivo{export\_best42.py}). Está preparado, pero requiere Linux o macOS y no
se ha ejecutado ni probado en Android.

\paragraph{Versión final del detector: YOLO26m en Ultralytics Platform.} La versión
final se entrenó y se publicó en Ultralytics Platform con YOLO26 en su tamaño
mediano (m), a partir de los pesos preentrenados que distribuye Ultralytics
(Jocher et~al., 2026; Ultralytics, s.~f.-c). YOLO26 es la generación más reciente
de la familia y sus autores la orientan a dispositivos de borde, con una
inferencia que no requiere supresión de no máximos (Sapkota et~al., 2025). El
conjunto de datos (\texttt{gym-machines-data}) se creó en el editor de la
plataforma, que permite anotar con polígonos y con la ayuda de un modelo de
segmentación (SAM) (Ultralytics, s.~f.-a): cada máquina se delimitó con un
polígono y se le asignó una de las 52 clases definidas para el proyecto. La
tarea de entrenamiento es de detección, de modo que el modelo predice cajas
delimitadoras y no contornos.

Según el gráfico de distribución de clases de la plataforma, el conjunto reúne
6\,681 ejemplos anotados (Tabla~\ref{tab:distribucion-clases}). La clase más
frecuente (\texttt{prensa-piernas}) aporta 615 ejemplos, el 9,2\,\% del total,
y cinco clases tienen cuatro ejemplos o menos (\texttt{rack-sentadilla},
\texttt{press-alto}, \texttt{press-asistida}, \texttt{extension-cuadriceps} y
\texttt{remo-alto}). Las 24 clases con 100 ejemplos o más concentran el
88,9\,\% de las anotaciones, y las seis con menos de diez apenas el 0,3\,\%. El
desbalance es fuerte (de 615 ejemplos a uno en la clase más rara) y, con tan
pocos ejemplos, esas clases difícilmente pueden aprenderse o evaluarse: su
resultado individual no es interpretable y las métricas globales se apoyan
sobre todo en las clases frecuentes. Además, algunas etiquetas parecen nombrar
la misma máquina con criterios distintos (por ejemplo,
\texttt{abductores-aductores} y \texttt{hip\_adduction\_machine}, o
\texttt{curl-biceps} y \texttt{bicep\_curl\_machine}), lo que puede repartir los
ejemplos de una máquina entre dos clases.

\begin{table}[H]
\centering
\caption{Distribución de los ejemplos anotados por clase en \texttt{gym-machines-data}}
\label{tab:distribucion-clases}
\small
\begin{tabular}{|l|C{1.8cm}|C{2.2cm}|C{2.6cm}|}
\hline
\rowcolor{encabezado} \textbf{Ejemplos por clase} & \textbf{Clases} &
\textbf{Ejemplos} & \textbf{\% del total} \\
\hline
100 o más & 24 & 5\,939 & 88,9\,\% \\ \hline
De 30 a 99 & 11 & 473 & 7,1\,\% \\ \hline
De 10 a 29 & 11 & 248 & 3,7\,\% \\ \hline
Menos de 10 & 6 & 21 & 0,3\,\% \\ \hline
\textbf{Total} & \textbf{52} & \textbf{6\,681} & \textbf{100\,\%} \\
\hline
\end{tabular}
\fuente{Elaboración propia con los conteos por clase del gráfico de
distribución de clases de Ultralytics Platform.}
\end{table}

El entrenamiento duró 80 épocas, con imágenes de 640$\times$640 píxeles, lotes
de 32, una GPU, semilla 42 en modo determinista y precisión mixta automática. El
optimizador se dejó en automático y se mantuvieron los aumentos de datos por
defecto de Ultralytics (mosaico, que se desactiva en las últimas 10 épocas,
volteo horizontal, variaciones de color y borrado aleatorio). No se configuró
parada temprana (paciencia igual a 0). La corrida se registró como
\texttt{version-5} del proyecto \texttt{detect-machine}. La
Tabla~\ref{tab:metricas-detector-yolo26} resume las métricas de la última época,
calculadas sobre la partición de validación.

\begin{table}[H]
\centering
\caption{Métricas del detector final (YOLO26m) en la partición de validación}
\label{tab:metricas-detector-yolo26}
\small
\begin{tabular}{|l|C{2.6cm}|}
\hline
\rowcolor{encabezado} \textbf{Métrica} & \textbf{Valor} \\
\hline
Precisión & 57,7\,\% \\ \hline
Exhaustividad (\emph{recall}) & 45,0\,\% \\ \hline
mAP50 & 48,5\,\% \\ \hline
mAP50--95 & 44,8\,\% \\ \hline
Pérdida de caja (entrenamiento / validación) & 0,243 / 0,612 \\ \hline
Pérdida de clasificación (entrenamiento / validación) & 0,286 / 2,690 \\
\hline
\end{tabular}
\fuente{Registro de entrenamiento de Ultralytics Platform (corrida
\texttt{version-5}, época 80). Definición de las métricas en Padilla et~al.
(2021). La validación no es una partición de prueba independiente.}
\end{table}

Estas cifras son las del propio entrenamiento y deben leerse con cautela. La
partición de validación es la que Ultralytics usa también para elegir el mejor
punto de control, de modo que tiende a favorecer al modelo; no se evaluó con una
partición de prueba independiente ni con fotografías de un gimnasio real. Además,
la pérdida de clasificación en validación (2,69) es unas diez veces la de
entrenamiento (0,29), y la de caja (0,61) más del doble (0,24), lo que apunta a
sobreajuste: el modelo se ajusta bien a las imágenes que ya vio y generaliza
menos a las nuevas. Por todo ello las cifras muestran que el procedimiento
funciona de extremo a extremo, no que el detector sea fiable en un gimnasio.

\paragraph{Dispositivo utilizado para las pruebas}

\paragraph{Dispositivo de prueba.} El teléfono de prueba es un Infinix HOT 50
Pro+ (modelo X6880), un equipo de gama media cuyas características resume la
Tabla~\ref{tab:dispositivo-prueba}. Según la lista de dispositivos compatibles
de Google, este modelo admite ARCore y la API de profundidad (Google for
Developers, s.~f.-a), que es lo que necesita el módulo CoachAR. Con un único
equipo no puede afirmarse que la aplicación se comporte igual en otros
teléfonos, con cámaras, memoria o soporte de ARCore distintos.

\begin{table}[H]
\centering
\caption{Características del dispositivo de prueba}
\label{tab:dispositivo-prueba}
\small
\begin{tabular}{|l|l|}
\hline
\rowcolor{encabezado} \textbf{Característica} & \textbf{Valor} \\
\hline
Modelo & Infinix HOT 50 Pro+ (X6880) \\ \hline
Procesador & MediaTek Helio G100 \\ \hline
Memoria RAM & 8\,GB (más 5\,GB de RAM extendida) \\ \hline
Almacenamiento & 256\,GB \\ \hline
Cámara trasera & 50\,MP + 2\,MP \\ \hline
Cámara frontal & 13\,MP \\ \hline
Pantalla & 1\,080 $\times$ 2\,436 píxeles \\ \hline
Batería & 5\,000\,mAh \\ \hline
Sistema operativo & XOS 15.1.2 (basado en Android 15) \\
\hline
\end{tabular}
\fuente{Pantalla «Mi teléfono» de los ajustes del dispositivo. Compatibilidad
con ARCore: Google for Developers (s.~f.-a).}
\end{table}





\paragraph{Verificación de la realidad aumentada}

La aplicación se probó en un dispositivo físico y cumplió los criterios de
verificación de la realidad aumentada. El modelo reconoció la máquina y el
módulo \archivo{CoachAR} mostró su información mediante ARCore. El seguimiento
espacial, la consulta al servicio de reconocimiento y la visualización de la
ficha funcionaron correctamente de forma integrada. Los casos CP-42 a CP-44
del Anexo~\ref{anx:casos} se completaron satisfactoriamente.

\paragraph{Evaluación del recomendador de rutinas}

\paragraph{Resultados.} En el conjunto de prueba el modelo acertó la plantilla
en el 99,8\,\% de los casos, el esquema de carga en el 99,97\,\% y ambas a la
vez en el 99,8\,\% (Tabla~\ref{tab:metricas-rutinas}). Como comprobación
adicional se evaluó el modelo sobre una cuadrícula exhaustiva de 28\,512
perfiles que recorre todas las combinaciones de objetivo, nivel, género, días,
edades y valores de IMC alrededor de los umbrales: coincidió con las reglas en
27\,689 (97,1\,\%). Los desacuerdos se concentran justo en los umbrales, por
ejemplo con 59 años (un año antes del cambio de esquema) o con IMC de 29,9
(justo antes del énfasis cardio), donde un modelo continuo suaviza la frontera
que las reglas trazan de forma abrupta.

\begin{table}[H]
\centering
\caption{Desempeño del recomendador de rutinas sobre datos sintéticos}
\label{tab:metricas-rutinas}
\small
\begin{tabular}{|l|C{3.4cm}|}
\hline
\rowcolor{encabezado} \textbf{Medida} & \textbf{Resultado} \\
\hline
Exactitud de la plantilla semanal (prueba, 9\,000 perfiles) & 99,8\,\% \\ \hline
Exactitud del esquema de carga (prueba) & 99,97\,\% \\ \hline
Exactitud de ambas clases (prueba) & 99,8\,\% \\ \hline
Coincidencia con las reglas en la cuadrícula exhaustiva (28\,512 perfiles) &
97,1\,\% \\ \hline
Coincidencia del archivo \texttt{.tflite} con el modelo de Keras & 100\,\% \\
\hline
\end{tabular}
\fuente{Archivos \archivo{ml/results/metrics.json} y
\archivo{ml/results/model\_metadata.json}.}
\end{table}

\paragraph{Comprobación del reporte PDF}
Según el registro consignado en el documento, el 3 de noviembre de 2026 se
generó un reporte de 30 días en el Infinix HOT 50 Pro+. El archivo incluyó
perfil, peso e IMC, entrenamientos y alimentación y se abrió en el visor del
teléfono. La Figura~\ref{fig:reporte-pdf} presenta datos de demostración.
Esta comprobación corresponde al reporte y no valida el módulo de ARCore.

\subsection{Pruebas funcionales en dispositivos}

Las pruebas funcionales se realizan con la técnica de caja negra sobre el
\emph{development build} instalado en un dispositivo Android, siguiendo los
criterios de aceptación definidos en cada iteración. Se definieron 44 casos de
prueba que cubren los requerimientos funcionales RF1 a RF6; la matriz completa,
con los pasos, el resultado esperado y el resultado obtenido, se encuentra en
el Anexo~\ref{anx:casos}. La Tabla~\ref{tab:resumen-funcionales} resume su
distribución.

\begin{table}[H]
\centering
\caption{Distribución de los casos de prueba funcionales}
\label{tab:resumen-funcionales}
\begin{tabular}{|L{6cm}|C{2.5cm}|C{3.4cm}|}
\hline
\rowcolor{encabezado} \textbf{Módulo} & \textbf{Casos} & \textbf{Requerimientos} \\
\hline
Autenticación, onboarding y perfil & 9 & RF1.1 a RF1.5 \\ \hline
Guía de máquinas y ejercicios & 6 & RF2.1 a RF2.3 \\ \hline
Escáner con reconocimiento y realidad aumentada & 7 & RF2.4 \\ \hline
Rutinas y planificación & 7 & RF3.1 a RF3.3 \\ \hline
Nutrición (alimentos e IMC) & 7 & RF4.1, RF5.1 a RF5.3 \\ \hline
Progreso, entrenamientos y reportes & 8 & RF4.2, RF6.1, RF6.2 \\
\hline
\textbf{Total} & \textbf{44} & \\
\hline
\end{tabular}
\end{table}

\subsection{Pruebas funcionales con usuarios}

\subsubsection{Pruebas de usabilidad}

Las pruebas de usabilidad permitieron comprobar la realización de las tareas
principales de la aplicación (Nielsen, 1993). Las seis tareas evaluadas se
completaron correctamente y cumplieron sus criterios de éxito. La
Tabla~\ref{tab:tareas-usabilidad} presenta los resultados obtenidos.

\begin{table}[H]
\centering
\caption{Resultados de las pruebas de usabilidad}
\label{tab:tareas-usabilidad}
\small
\begin{tabularx}{\textwidth}{|L{0.7cm}|Y|Y|L{1.6cm}|}
\hline
\rowcolor{encabezado} \textbf{ID} & \textbf{Tarea} & \textbf{Criterio de éxito} & \textbf{Resultado} \\
\hline
T1 & Registrarse y completar el \emph{onboarding}. & Llega a la pantalla de
inicio. & Cumple \\ \hline
T2 & Encontrar la ficha de la prensa de piernas. & Abre la ficha desde el
catálogo. & Cumple \\ \hline
T3 & Escanear una máquina y abrir su ficha completa. & Abre la ficha desde la
superposición. & Cumple \\ \hline
T4 & Generar una rutina y marcar como completado el día de hoy. & La rutina
queda activa y el día marcado. & Cumple \\ \hline
T5 & Registrar 150~g de arroz en el diario. & El alimento aparece en la lista
del día. & Cumple \\ \hline
T6 & Registrar su peso y ver el gráfico de evolución. & Ve el nuevo punto en
el gráfico. & Cumple \\
\hline
\end{tabularx}
\end{table}

\subsection{Pruebas de aceptación por objetivo}

La Tabla~\ref{tab:trazabilidad} relaciona cada objetivo específico con los
requerimientos que lo implementan y con las pruebas que lo verifican. Esta
trazabilidad es la base de la evaluación del cumplimiento de objetivos que
realiza el tribunal (Artículo~64 del Reglamento de Titulación).

\begin{table}[H]
\centering
\caption{Trazabilidad entre objetivos, requerimientos y pruebas}
\label{tab:trazabilidad}
\small
\begin{tabularx}{\textwidth}{|L{1.1cm}|Y|L{2.2cm}|L{3.4cm}|}
\hline
\rowcolor{encabezado} \textbf{Obj.} & \textbf{Funcionalidad entregada} & \textbf{Req.} &
\textbf{Pruebas} \\
\hline
OE1 & Guía de 35 máquinas con búsqueda, filtros, fichas y escáner con
reconocimiento asistido (modelo local o detector en la nube) y ficha anclada
con ARCore en los teléfonos compatibles. & RF2.1, RF2.2, RF2.4 &
CP-10 a CP-13, CP-15 a CP-19, CP-42 a CP-44, sección~\ref{sec:pruebas-ia} \\
\hline
OE2 & Ejercicios con animación, músculo objetivo y dificultad. & RF2.3 &
CP-14 \\ \hline
OE3 & \emph{Onboarding} de objetivo, nivel y datos físicos; rutinas
generadas por un recomendador de TensorFlow Lite con respaldo de reglas. &
RF1.4, RF3.1 & CP-06 a CP-08, CP-20, CP-21,
Tabla~\ref{tab:prueba-generador}, secciones~\ref{sec:pruebas-unitarias}
y~\ref{sec:pruebas-ia} \\ \hline
OE4 & Planificador de rutinas por día, rutina activa y agregado desde la
ficha. & RF3.2, RF3.3 & CP-22 a CP-26 \\ \hline
OE5 & Diario de alimentación con calorías y macronutrientes, metas
personalizadas; calculadora de IMC con clasificación e historial. & RF4.1, RF5 &
CP-27 a CP-33, Tablas~\ref{tab:prueba-imc} y~\ref{tab:prueba-otras},
sección~\ref{sec:pruebas-unitarias} \\ \hline
OE6 & Registro de peso, estatura e IMC. & RF4.2 & CP-34, CP-35 \\ \hline
OE7 & Reporte con gráfico, variación y categoría actual; reporte en PDF por
período; seguimiento de entrenamientos con racha. & RF6 & CP-36 a CP-41,
Tabla~\ref{tab:prueba-otras} \\
\hline
\end{tabularx}
\end{table}

\subsection{Cumplimiento del alcance y desviaciones}

La Tabla~\ref{tab:desviaciones} contrasta lo planificado en los
Capítulos~\ref{cap:introduccion} y~\ref{cap:area-aplicacion} con lo que
efectivamente se entrega. Se incluye para que el alcance del producto se
evalúe con transparencia.

\begin{table}[H]
\centering
\caption{Contraste entre lo planificado y lo implementado}
\label{tab:desviaciones}
\small
\begin{tabularx}{\textwidth}{|L{3.1cm}|L{2.4cm}|Y|}
\hline
\rowcolor{encabezado} \textbf{Elemento planificado} & \textbf{Estado} & \textbf{Observación} \\
\hline
Rutinas personalizadas con IA mediante TensorFlow Lite & Completo, con
salvedad & Un modelo entrenado decide la estructura de la rutina en el
dispositivo y las reglas actúan como respaldo. Se entrenó con perfiles
sintéticos etiquetados por las reglas, por lo que imita el criterio experto y
no se validó con resultados de usuarios reales. \\ \hline
Guía y escaneo de máquinas con realidad aumentada & Parcial & Implementado:
reconocimiento local (TensorFlow Lite, 399 fotos de catálogo) o en la nube
(detector YOLO26m con cajas) y anclaje de la ficha con ARCore en un módulo
nativo propio, con cámara convencional como alternativa. Verificado en un
dispositivo físico: el seguimiento espacial y la consulta al servicio de
reconocimiento funcionaron correctamente durante las pruebas. El anclaje se fija a un punto de una superficie,
no a la geometría de la máquina, y el reconocimiento ocurre por captura, no de
forma continua. El acierto es modesto (38,5\,\% del modelo local en campo
simulado; detector con mAP50 de 48,5\,\% en validación): solo
se acepta automáticamente la máquina con confianza $\geq 0{,}8$ y en los demás
casos el usuario confirma. \\ \hline
Registro nutricional (solo IMC según el alcance) & Ampliado & Se añadió un
diario de alimentación con USDA y metas diarias personalizadas con la ecuación
de Mifflin-St Jeor. \\ \hline
Progreso y reportes & Completo & Gráfico de peso e IMC, variación,
categoría, seguimiento de entrenamientos con racha y reporte en PDF por período, verificado en el dispositivo de prueba (CP-41). \\ \hline
Sincronización en tiempo real & Completo & Escuchas \emph{onSnapshot}; no se
configuró la persistencia de datos sin conexión. \\ \hline
Android 7.0 (API 24) o superior & Verificado en Android 15 & La aplicación
funcionó correctamente en el dispositivo de prueba con Android 15. Esta
comprobación corresponde a esa versión; la compatibilidad con Android 7.0 no
se verificó directamente. \\ \hline
Componente de servidor (no previsto en la planificación inicial) & Agregado &
Las funciones \archivo{buscarAlimentos}, \archivo{guardarRutina} y
\archivo{reconocerMaquina} (Cloud Functions) protegen la clave de USDA, validan
la creación de rutinas y ocultan la clave del servicio de reconocimiento;
añaden dependencia de la infraestructura de Firebase y de Ultralytics. \\ \hline
Dispositivos externos (\emph{smartwatch}) & Fuera de alcance & Excluidos por el
alcance del proyecto. \\
\hline
\end{tabularx}
\end{table}
